\subsubsection{Layer 5: Gauss Digamma = Bloch Decomposition}

The Gauss digamma formula
\[
  \psi(p/q) = -\gamma - \ln(2q) - \frac{\pi}{2}\cot\left(\frac{\pi p}{q}\right)
  + 2\sum_{r=1}^{\lfloor(q-1)/2\rfloor} \cos\left(\frac{2\pi rp}{q}\right)
  \ln\sin\left(\frac{\pi r}{q}\right)
\]
connects the logarithmic series (arising from the telescope) to the
cotangent sums that appear in the Vasyunin formula.

\begin{correspondence}[Bloch Theorem on the Prime Lattice]
The Gauss digamma formula is the \emph{Bloch decomposition} of the
digamma function on the cyclic group $\ZZ/q\ZZ$: the wavefunctions on
a periodic lattice with $q$ sites decompose into discrete Fourier
modes $e^{2\pi irp/q}$, and the digamma at the rational point $p/q$
is the sum over these modes weighted by $\ln\sin(\pi r/q)$---the
\emph{Bloch energies} of the lattice.

In solid-state physics, the Bloch theorem says that eigenstates of
a periodic Hamiltonian can be labelled by crystal momentum $k$,
and the energy $E(k)$ is a periodic function of $k$. Here,
$r$ plays the role of crystal momentum and $\ln\sin(\pi r/q)$
is the dispersion relation.

This is currently an axiom (\lean{gauss\_digamma\_formula} in
\lean{DigammaReflection.lean}), provable from Mathlib's digamma
infrastructure and the Fourier expansion of $\ln\Gamma$.
\end{correspondence}

\subsubsection{Layer 6: The Convergence Axiom as Ergodic Hypothesis}

The final axiom (\lean{partial\_integral\_tends\_to\_formula} in
\lean{ConvergenceAxioms.lean}) states that the partial row-sum
$\int_{1/(aM)}^1 \fract{1/(ax)}\fract{1/(bx)}\,dx$ converges
to the Vasyunin formula as $M \to \infty$.

\begin{principle}[The Ergodic Hypothesis]
The convergence axiom is the number-theoretic \emph{ergodic hypothesis}:
the time average (partial lattice sum as $M \to \infty$) equals the
ensemble average (closed-form Vasyunin formula).

The Cathedral proves one half of this hypothesis mechanically:
\begin{itemize}
  \item \textbf{Route A (proved):} The partial integral converges to
    \lean{gramIntegral}~$= \int_0^1 \fract{1/(ax)}\fract{1/(bx)}\,dx$
    because the tail $\int_0^{1/(aM)}$ is squeezed to zero.
    This is the \emph{thermodynamic limit}: as the lattice size
    $M \to \infty$, the finite-volume partition function converges
    to the infinite-volume limit.
  \item \textbf{Route B (axiom):} The partial integral also converges
    to the Vasyunin formula. This is the \emph{ergodic hypothesis proper}:
    the raw partition function agrees with the closed-form evaluation.
\end{itemize}
The \lean{tendsto\_nhds\_unique} lemma (uniqueness of limits in
Hausdorff spaces) then forces \lean{gramIntegral} = Vasyunin formula.
This is the statement that the system is ergodic: since both limits
exist and sequences in Hausdorff spaces have unique limits, the
thermodynamic limit must equal the ensemble average.
\end{principle}

\subsection{The Cotangent Sum as Scattering Phase}

The Vasyunin cotangent sum
\[
  V(a,b) = \sum_{m=1}^{a-1} \fract{mb/a} \cdot \cot(\pi m/a)
\]
encodes the interaction between modes $a$ and $b$ through the
\emph{modular structure} of their ratio $b/a$.

\begin{correspondence}[Scattering Phase Shift]
$V(a,b)$ is the \emph{scattering phase shift} between modes
$a$ and $b$ on the prime lattice. The fractional part
$\fract{mb/a}$ encodes the relative displacement (the modular
position of mode $b$ at lattice site $m$ of mode $a$), and the
cotangent $\cot(\pi m/a)$ is the Hilbert kernel---the boundary
value of the Cauchy kernel on the unit circle.

The off-diagonal Gram entry decomposes as:
\begin{align}
  G(a,b) &= \underbrace{\frac{\ln 2\pi - \gamma}{2}\left(\frac{1}{a}+\frac{1}{b}\right)}_{\text{vacuum (zero-point energy)}}
  + \underbrace{\frac{a-b}{2ab}\ln\frac{b}{a}}_{\text{Born approximation}}
  \nonumber \\
  &\quad - \underbrace{\frac{\pi}{2ab}\bigl(V(a,b) + V(b,a)\bigr)}_{\text{scattering phases}}
  - \underbrace{\frac{1}{ab}}_{\text{contact term}}
\end{align}
for coprime $a < b$.
Each term has a precise physical meaning:
\begin{itemize}
  \item \textbf{Vacuum}: The $\ln 2\pi - \gamma$ terms are the
    one-loop entropy and Euler--Mascheroni bare mass, independent of
    the specific modes.
  \item \textbf{Born approximation}: The $\ln(b/a)$ term is the
    tree-level scattering amplitude, depending only on the ratio of
    ``momenta'' $a$ and $b$.
  \item \textbf{Scattering phases}: $V(a,b) + V(b,a)$ encodes the
    full interaction via modular arithmetic---the ``non-perturbative''
    contribution.
  \item \textbf{Contact term}: $1/(ab)$ is the delta-function
    interaction at zero separation.
\end{itemize}
\end{correspondence}

The Dedekind reciprocity axiom (\lean{harmonicTileSum\_reciprocity}),
which states a reciprocity law for harmonic tile sums, is the number-theoretic
incarnation of \emph{time-reversal invariance}: the scattering amplitude
for $a \to b$ is related to $b \to a$ by a universal kinematic factor.

